% Chapter 1 of Deep Learning as an Experimental Science: seeds and the run distribution.
\newcommand{\seedsMS}[1]{\mathrm{MS}_{#1}}
\labchapter{Seeds}{The spread of a single method}{One run of each method is enough to rank them.}{Common practice in leaderboard tables, 2015--2019; audited for named-entity recognition by Reimers and Gurevych (2017) and for BERT fine-tuning by Dodge et al.\ (2020)}{Vanished: changing only the seed moves a score by as much as the gaps between the systems that single runs were used to rank.}{ch:seeds}
\begin{refsection}[chapters/01-seeds.bib]

\begin{chaptersummary}
Until about 2019 most tables in deep learning compared one trained network per method, and gaps of a few tenths of a point decided the state of the art. Re-running published systems with nothing changed but the random seed moved their scores by as much as those gaps and put one celebrated pair of named-entity taggers in the opposite order, and BERT's scores on small tasks turned out to depend on how many restarts had been tried. The method treats a training recipe as a distribution of runs: we estimate its mean and spread, split the spread into the parts due to initialization, data order, data split and the finite test set, compute what reporting the best of $k$ runs adds, and work out how many seeds a comparison needs. In the laboratory, two hundred runs of one small network spread over nineteen test items, and the best of ten runs gains almost a point of test error that fresh test items do not confirm.
\end{chaptersummary}

\section{The claim: tenths of a point on CoNLL-2003}\label{sec:seeds:claim}

In 2016 two groups published neural taggers for named-entity recognition, the task of finding the names of persons, organizations, locations and other entities in text, and measured them on the CoNLL-2003 shared task \parencite{TjongKimSangDeMeulder2003}: Reuters news stories with a fixed split into training, development and test sets, whose English test set holds 5,648 entity mentions, scored by F1.\sidenote{F1 is the harmonic mean of precision and recall over entity mentions; a mention counts as found when both its span and its type are right.} \Textcite{LampleEtAl2016} combined a bidirectional LSTM, character-level LSTMs and a conditional random field, and reported 90.94 F1; \textcite{MaHovy2016} built the character representations with a convolutional network instead and reported 91.21.

\begin{finding}[\textcite{MaHovy2016}, Table~5]\label{fd:seeds:ner}
On the English test set of CoNLL-2003, the BiLSTM-CNN-CRF tagger reached 91.21 F1, against 90.94 for the LSTM-CRF of \textcite{LampleEtAl2016} and 91.20 for the best earlier system, a joint model of entity recognition and entity linking. The five best entries of the table lie within 0.31 points of each other. Each entry is one number; neither Ma and Hovy nor Lample et al.\ report a spread over runs.
\end{finding}

Ma and Hovy reported improvements over earlier neural taggers, Lample et al.'s among them, and a new state of the art by 0.01 points over the joint model, a margin they called slight. The reading that such tables invited was simpler than either paper, and it is the claim of this chapter: the larger number belongs to the better system, so one run of each method is enough to rank them.

Practice was not uniform: \textcite{ChiuNichols2016} reported means and standard deviations over ten runs (91.62 and 0.33 for their best tagger), after excluding the 5--10\% of runs whose training failed, identified on training data,\sidenote{A failure rule computed on training data is part of the recipe and applies to every new run; a rule computed on test scores is selection on the test set.} and \textcite{PetersEtAl2018} reported ELMo's result as 92.22 with standard deviation 0.10 over five seeds. But ranking tables set such means beside single runs and best runs, and neighboring rows differed by about the standard deviations that the careful papers reported.

Two years later the question returned at a larger scale, on the GLUE benchmark, whose test labels sit on a server that scores submissions.

\begin{finding}[\textcite{DevlinEtAl2019}, Section~4.1 and Table~1]\label{fd:seeds:bert}
Fine-tuned BERT-Large outperformed all earlier systems on every GLUE task; on the smallest, RTE, with 2,500 training examples, its test accuracy was 70.1 against 56.0 for OpenAI GPT. Because fine-tuning was sometimes unstable on small datasets, the authors ran several restarts with different data shuffling and classifier initialization, and submitted the one that was best on the development set.
\end{finding}

The restarts were sound engineering, and the paper reports them: the selection used development data, and each BERT model went to the test server once.\sidenote{Footnote~9 of \textcite{DevlinEtAl2019}. In their ablations and their CoNLL-2003 experiments the same authors report averages over five restarts.} But they change what the reported number measures: not one run of a recipe, but the best of several. And the GLUE leaderboard, like the CoNLL tables, ranked systems by one number each.

The practice was reasonable. For a linear model trained to the optimum of a strictly convex objective, such as a regularized conditional random field, a second run reproduces the first and one run is the method; the habit carried over to networks, whose training is a random experiment. Runs were expensive (about six hours on one processor for a tagger of \textcite{ChiuNichols2016}), and practice was improving: in a survey of authors of NeurIPS 2019 and ICLR 2020 papers, a typical study reported about four results per model and task \parencite{BouthillierVaroquaux2020}. The original authors also did what this chapter relies on: Ma and Hovy and Lample et al.\ released the code that made the re-runs possible, and Devlin et al.\ kept the test set out of their selection.

\section{The experiment: one run per arm}\label{sec:seeds:experiment}

Stripped of its details, each comparison in \cref{fd:seeds:ner} has the following design.

\begin{experiment}[a leaderboard comparison, CoNLL-2003, 2016]
\field{question}{Does tagger A recognize named entities better than tagger B?}
\field{treatment}{the system: architecture, features, embeddings and training recipe; the unit is a training run}
\field{levels}{the published systems, e.g.\ an LSTM-CRF with character LSTMs and a BiLSTM-CNN-CRF}
\field{tuned}{by each team, on the development set, with its own budget}
\field{held fixed}{the English data of CoNLL-2003, its standard split, and the F1 scorer}
\field{varied}{seeds of initialization, data order and dropout; software and hardware; none recorded}
\field{response}{F1 on the 3,684 sentences of the test set}
\field{replicates}{one run per system, or the best of several on the development set}
\field{analysis}{compare the two numbers; the larger belongs to the better system}
\end{experiment}

The GLUE submissions of 2018 and 2019 had the same design, with the best of several restarts in place of one run for some systems.

What could this design show? If two runs of one tagger always agreed to within a few hundredths of a point, then the larger number would identify the better system whenever the gap exceeded those hundredths. The design cannot tell: with one run per arm nothing estimates how far a second run would land, so the comparison has no error bar and rests on an unmeasured number, the standard deviation of a tagger's F1 over seeds. Two habits made the gaps look firmer. A significance test over the test sentences, such as a paired bootstrap, treats the two trained networks as fixed and measures only the noise of the test set. And selecting the best run on the development set made the reported number a best of several runs, a statistic with a different expectation (\cref{prop:seeds:bestof}).

\section{The post-mortem: the seed alone}\label{sec:seeds:postmortem}

\paragraph{Re-running the taggers.} \Textcite{ReimersGurevych2017} re-ran both taggers with the authors' public code, changing nothing but the seed: Ma and Hovy's 86 times and Lample et al.'s 41 times. Ma and Hovy's tagger scored between 89.99 and 91.00 F1 (median 90.64, standard deviation 0.24), Lample et al.'s between 90.19 and 91.14 (median 90.81, standard deviation 0.18). A Kolmogorov--Smirnov test separated the two distributions ($p<0.01$), with Lample et al.'s quartiles above Ma and Hovy's, the opposite of the published order, and the published 91.21 lay above all 86 re-runs. Selection on the development set did not rescue a single run: across Lample et al.'s runs the rank correlation between development and test F1 was 0.23, and the run with the best development score scored 90.31 on the test set.\sidenote{The audit was part of a study of 50,000 LSTM taggers on five tasks, which looked for architectures that are both accurate and stable \parencite{ReimersGurevych2017}.} By \cref{prop:seeds:tworuns} below, two runs of Ma and Hovy's tagger differ by more than the published gap of 0.27 with probability $2\Phi(-0.27/(0.24\sqrt2))\approx0.43$, and, treating the medians as means, one run of each puts Ma and Hovy's tagger ahead, the published order, with probability $\Phi(-0.17/0.30)\approx0.28$.

\paragraph{Two runs of every configuration.} The same study trained 1,830 sampled tagger configurations twice each. On CoNLL-2003 the two runs of a configuration differed by a median of 0.38 points, with 95th percentile 1.08 and maximum 2.59. Under normality these quantiles are $0.954\sigma$ and $2.77\sigma$ (\cref{prop:seeds:tworuns}), so both point to a standard deviation near 0.4 points, more than the 0.31 points that separate the five best systems of \cref{fd:seeds:ner}. On another entity task one configuration scored 82.5 with one seed and 74.3 with another.

\paragraph{A test over test sentences does not see the seed.} In an unrefereed follow-up, \textcite{ReimersGurevych2018} trained one tagger configuration 100,000 times per task and compared it with itself by a bootstrap test over the test sentences at level 0.05. The test declared two runs of the configuration different in 31\% of comparisons on English CoNLL-2003 (11--34\% on seven tasks), and still in 22\% when each side was the best of 50 runs on development data. As the authors note, the test is not at fault: it asks whether two fixed networks differ, and two runs are two networks. The error is the step from different networks to different methods, which needs the variation between runs that a test over sentences never sees (\cref{prop:seeds:noise}).

\paragraph{Fine-tuning BERT.} \Textcite{DodgeEtAl2020}, in a preprint, fine-tuned BERT-Large on four GLUE tasks with fixed hyperparameters, varying only the seed of the classifier layer (the only new weights) and the seed of the data order, in all $25\times25$ combinations on MRPC, RTE and CoLA and $15\times15$ on SST-2. The best validation score beat the published validation scores of BERT by about seven points on RTE and CoLA (on RTE, 77.3 against 70.0) and, on MRPC, the scores of the later XLNet, RoBERTa and ALBERT. The two seeds contributed comparably to the variance, and on the small tasks some runs diverged to the accuracy of predicting the majority class, making the distribution of runs bimodal. A maximum over 625 runs carries the selection bonus of \cref{prop:seeds:bestof}; Dodge et al.\ also plot how the expected best grows with the number of runs, and their point is that the gap between a published BERT and a later model is in part a gap in restarts (\cref{ch:budgets}).

\paragraph{Ten thousand seeds.} \Textcite{Picard2021}, also a preprint, trained a small residual network on CIFAR-10 for 30 seconds with each of $10^4$ seeds (mean accuracy 90.02, standard deviation 0.23, range 89.01 to 90.83) and fine-tuned three pretrained ImageNet models with 50 seeds each (standard deviations 0.08--0.11, ranges 0.37--0.46 points).\sidenote{The seed of the ornament at the foot of each chapter opening in this book, 3407, is the one in Picard's title.} These extremes are what a normal distribution predicts: the expected range of $k$ normal draws, $2\sigma e_k$ (\cref{prop:seeds:bestof}), is 1.77 for $\sigma=0.23$ and $k=10^4$, and 0.45 for $\sigma=0.10$ and $k=50$. Picard's black swans are the tails of an ordinary distribution sampled many times.

\paragraph{More than the seed.} \Textcite{BouthillierEtAl2021} varied the sources of randomness one at a time, 200 runs per source, in five case studies from BERT on GLUE to VGG on CIFAR-10. Resampling the data, which changes both the training and the test set, was the largest source, and a binomial model of the finite test set accounted for it; initialization was generally less than half of it and on a par with data order, and repeating the hyperparameter search added about as much as initialization. The authors recommend randomizing as many sources as possible.

\paragraph{Why the failed runs failed.} The failures behind the restarts were commonly attributed to small datasets. \Textcite{MosbachEtAl2021} and \textcite{ZhangEtAl2021} traced them to the optimization recipe: the version of Adam used to fine-tune BERT omitted bias correction, which matters most in the first steps, and three epochs on a small dataset are few steps. With bias correction and longer training the failures disappeared, and a reduced training set caused almost none once the number of steps was held fixed (\cref{ch:mechanisms}). On RTE over 25 seeds, Devlin et al.'s recipe gave a mean development accuracy of 50.9 (standard deviation 4.5, maximum 67.5), Mosbach et al.'s 67.3 (2.7, 71.1): the best run moved by 3.6 points, the mean by 16.4.

\begin{verdict}{vanished}
Under a protocol that runs each system many times, the gaps that single runs were used to rank fall inside the spread of one system's runs: the F1 of a tagger moves with its seed by a standard deviation of 0.2--0.4 points, against published gaps of 0.01--0.3, and the one pair of taggers that was re-run many times came out in the opposite order. What survives is a condition, not a ranking: one run per arm ranks two recipes reliably only when their gap is several seed standard deviations, and that standard deviation has to be measured.
\end{verdict}

\section{The method: the run distribution and its two kinds of noise}\label{sec:seeds:distribution}

\begin{definition}[Run distribution]\label{def:seeds:rundist}
A \term{training recipe}\index{training recipe} fixes everything about a training run that the experimenter chooses: the data and the rule that splits them, the model, the optimizer, the schedule, the hyperparameters and the code. Together with a rule for drawing its seeds, a recipe defines a probability distribution of outcomes, its \term{run distribution}\index{run distribution}. A run is one draw from it, and its \term{seeds}\index{seed} are the random inputs that select the draw.
\end{definition}

A claim about a recipe is a claim about its run distribution, usually about its mean $\mu$. The kit \script{labkit.py} keeps four sources of randomness apart: the initial weights (\texttt{init}), the order of the minibatches (\texttt{order}), the split into training and test sets (\texttt{split}), and the noise of label corruption, dropout and augmentation (\texttt{noise}). Nondeterministic kernels make even a run with every seed fixed a fresh draw (\cref{ch:instability}), and the hyperparameter search has randomness of its own (\cref{ch:budgets}). Which sources to randomize depends on the claim: a claim about a recipe on a fixed benchmark split randomizes initialization, order and noise, and a claim about a method on a task randomizes the split as well.

\paragraph{Summaries.} From $n$ runs with outcomes $Y_1,\dots,Y_n$, the sample mean $\bar Y$ estimates $\mu$, the sample standard deviation $S$ estimates the spread $\sigma$ of single runs, and the \term{standard error}\index{standard error} $\SE=S/\sqrt n$ the spread of $\bar Y$: $S$ says how far the next run may land from the mean, $\SE$ how far the mean of these runs may be from $\mu$. For normal runs $\bar Y\pm t_{n-1,0.975}\SE$ is a 95\% confidence interval for $\mu$ (\cref{thm:toolkit:tinterval}), and a bootstrap over runs gives an approximate one without normality (\cref{thm:toolkit:bootstrap}).\sidenote{The $t$ distribution was published by Gosset under the pen name Student \parencite{Student1908}, for the mean of a handful of measurements: the regime of a training experiment with five seeds.} Plotted against the number $j$ of runs so far, the \term{running mean}\index{running mean} and its interval show the estimate settling, the interval narrowing like $1/\sqrt j$. The foot of every chapter opening in this book draws this plot for sixteen standard normal draws, showing as many draws as the chapter's number; this chapter's opening shows one run and no interval, which is all that one run supports.

\begin{proposition}[One run of each]\label{prop:seeds:tworuns}
Let $Y_1\sim\cN(\mu_1,\sigma_1^2)$ and $Y_2\sim\cN(\mu_2,\sigma_2^2)$ be the outcomes of one run of each of two recipes, independent, with larger outcomes better and $\Delta=\mu_1-\mu_2\ge0$.
\begin{enumerate}
\item One run of each ranks the recipes in the wrong order with probability $\Prob(Y_1<Y_2)=\Phi\bigl(-\Delta/\sqrt{\sigma_1^2+\sigma_2^2}\bigr)$.
\item If the two recipes are the same, with $\sigma_1=\sigma_2=\sigma>0$, then $\Prob(\abs{Y_1-Y_2}>\delta)=2\Phi\bigl(-\delta/(\sigma\sqrt2)\bigr)$ for every $\delta\ge0$, and the median and the 95th percentile of $\abs{Y_1-Y_2}$ are $\sqrt2\,z_{0.75}\,\sigma\approx0.954\,\sigma$ and $\sqrt2\,z_{0.975}\,\sigma\approx2.77\,\sigma$.
\end{enumerate}
\end{proposition}

\begin{proof}
The difference $Y_1-Y_2$ is normal with mean $\Delta$ and variance $\sigma_1^2+\sigma_2^2$, so $\Prob(Y_1<Y_2)=\Prob(Y_1-Y_2<0)=\Phi(-\Delta/\sqrt{\sigma_1^2+\sigma_2^2})$. In part~(ii) the variable $Z=(Y_1-Y_2)/(\sigma\sqrt2)$ is standard normal, so $\Prob(\abs Z>x)=2\Phi(-x)$ by symmetry, which gives the tail probability with $x=\delta/(\sigma\sqrt2)$. The $q$-quantile of $\abs Z$ is $z_{(1+q)/2}$, and the quantiles of $\abs{Y_1-Y_2}$ are $\sigma\sqrt2$ times those of $\abs Z$.
\end{proof}

With $\delta=\sigma$ the tail probability is $2\Phi(-1/\sqrt2)\approx0.48$: two runs of one recipe differ by more than a standard deviation about half the time. With $\Delta=\sigma$, one run of each ranks two recipes wrongly with probability $\Phi(-1/\sqrt2)\approx0.24$. A comparison by single runs is informative only when the gap is several times $\sigma$. The outcome of a classifier is moreover a measurement on a finite test set, which adds a second source of noise.

\begin{proposition}[Two sources of noise in a measured error]\label{prop:seeds:noise}
Let a run $R$ be drawn from the run distribution and test items $\mathbf Z=(Z_1,\dots,Z_m)$ be drawn from the data distribution, independently of each other and of $R$. Write $e_R(z)\in\{0,1\}$ for the indicator that the network of run $R$ errs on item $z$, $p(R)=\E[e_R(Z_1)\mid R]$ for its population error, $q(z)=\E\,e_R(z)$ for the probability that a fresh run errs on item $z$, and $\bar p=\E\,p(R)=\E\,q(Z_1)$. The measured error $\hat p=\frac1m\sum_{i=1}^me_R(Z_i)$ satisfies the following.
\begin{enumerate}
\item $\E\hat p=\bar p$ and $\Var\hat p=\Var p(R)+\E[p(R)(1-p(R))]/m$.
\item For a test set held fixed while the runs vary, the variance over runs has mean
\begin{equation}\label{eq:seeds:churn}
\E\bigl[\Var(\hat p\mid\mathbf Z)\bigr]=\frac{m-1}{m}\Var p(R)+\frac1m\E\bigl[q(Z_1)\bigl(1-q(Z_1)\bigr)\bigr].
\end{equation}
\item $\bar p(1-\bar p)=\E[q(Z_1)(1-q(Z_1))]+\Var q(Z_1)$; in particular $\E[q(Z_1)(1-q(Z_1))]\le\bar p(1-\bar p)$, with equality if and only if $q(Z_1)=\bar p$ almost surely.
\end{enumerate}
\end{proposition}

\begin{proof}
The plan is to condition on the run for part~(i), to derive part~(iii) by algebra, and to obtain part~(ii) from both by conditioning on the test items.

\emph{Part (i).} Given $R$, the indicators $e_R(Z_1),\dots,e_R(Z_m)$ are independent Bernoulli variables with mean $p(R)$, because the items are independent of each other and of $R$. Hence $\E[\hat p\mid R]=p(R)$ and $\Var(\hat p\mid R)=p(R)(1-p(R))/m$. Taking expectations gives $\E\hat p=\bar p$, and the law of total variance gives the formula for $\Var\hat p$.

\emph{Part (iii).} Since $q(1-q)=q-q^2$, we have $\E[q(Z_1)(1-q(Z_1))]=\bar p-\E\,q(Z_1)^2=\bar p-\bar p^2-\Var q(Z_1)$, which is the identity; the inequality and its case of equality follow because $\Var q(Z_1)\ge0$, with equality exactly when $q(Z_1)$ is almost surely constant.

\emph{Part (ii).} Given the items, $\hat p$ has mean $\frac1m\sum_iq(Z_i)$, an average of $m$ independent copies of $q(Z_1)$, whose variance is $\Var q(Z_1)/m$. The law of total variance, now conditioning on the items, gives $\E[\Var(\hat p\mid\mathbf Z)]=\Var\hat p-\Var q(Z_1)/m$. Inserting part~(i), in which $\E[p(R)(1-p(R))]=\bar p-\E\,p(R)^2=\bar p(1-\bar p)-\Var p(R)$, gives
\[
\E\bigl[\Var(\hat p\mid\mathbf Z)\bigr]=\Var p(R)+\frac{\bar p(1-\bar p)-\Var p(R)-\Var q(Z_1)}{m}.
\]
By part~(iii), $\bar p(1-\bar p)-\Var q(Z_1)=\E[q(Z_1)(1-q(Z_1))]$, and collecting the two terms in $\Var p(R)$ gives \eqref{eq:seeds:churn}.
\end{proof}

Part~(i) is the noise of the test set: even a fixed network measures with standard error $\sqrt{\bar p(1-\bar p)/m}$, about 0.7 points for an error of 3\% on 539 items. Part~(ii) says what a spread over seeds on a fixed benchmark measures: the variance of the population error plus the \term{item churn}\index{item churn} $\E[q(1-q)]/m$, the variance due to runs that err on different items. By part~(iii) the churn lies between zero, when every run errs on the same items, and the full binomial variance, when every item is equally hard; the rest of the binomial variance, $\Var q(Z_1)/m$, is \term{item difficulty}\index{item difficulty}, which moves all runs together when the test set changes. A seed standard deviation measured on a small test set can therefore be mostly a property of the test set, as \cref{lab:seeds:bestof} finds.


\section{The method: selection, components and the number of seeds}\label{sec:seeds:components}

\paragraph{The best of $k$.}\index{best of $k$} A reported number is often the best of several runs: of a few restarts, of the seeds tried during development, of a few submissions. Its expectation is not the mean of the recipe.

\begin{proposition}[The best of $k$ runs]\label{prop:seeds:bestof}
Let $k\ge1$, and let $e_k=\E\max_{i\le k}Z_i$ for independent standard normals $Z_1,\dots,Z_k$.
\begin{enumerate}
\item If $Y_1,\dots,Y_k\iid\cN(\mu,\sigma^2)$, then $\E\max_iY_i=\mu+\sigma e_k$ and $\E\min_iY_i=\mu-\sigma e_k$. Moreover $e_1=0$, the sequence $e_k$ increases strictly, and $e_k\le\sqrt{2\ln k}$.
\item If $(D_1,T_1),\dots,(D_k,T_k)$ are independent pairs, each bivariate normal with means $\mu_D,\mu_T$, standard deviations $\sigma_D,\sigma_T>0$ and correlation $\rho$, and $\iota$ is the index of the largest $D_i$, then $\E T_\iota=\mu_T+\rho\,\sigma_Te_k$.
\end{enumerate}
\end{proposition}

\begin{proof}
The plan is to reduce part~(i) to standard normals and bound $e_k$ through the moment generating function, and then to derive part~(ii) from part~(i) by conditioning on the $D_i$.

\emph{Part (i).} Write $Y_i=\mu+\sigma Z_i$ with $Z_i$ independent standard normals. Since $\sigma\ge0$, $\max_iY_i=\mu+\sigma\max_iZ_i$, which gives the first formula; and $\min_iY_i=\mu-\sigma\max_i(-Z_i)$, where the $-Z_i$ are again independent standard normals, which gives the second. For $k=1$, $e_1=\E Z_1=0$. The maximum of $Z_1,\dots,Z_{k+1}$ is at least that of $Z_1,\dots,Z_k$, and strictly larger on the event that $Z_{k+1}$ exceeds the others, which has probability $1/(k+1)>0$; so $e_{k+1}>e_k$. For the bound, let $k\ge2$ and $t>0$. Jensen's inequality for the convex exponential, followed by bounding the largest of $k$ positive numbers by their sum, gives
\[
e^{te_k}\le\E\,e^{t\max_iZ_i}\le\sum_{i=1}^k\E\,e^{tZ_i}=k\,e^{t^2/2}.
\]
Taking logarithms, $e_k\le(\ln k)/t+t/2$, and the choice $t=\sqrt{2\ln k}$, which minimizes the right side, gives $e_k\le\sqrt{2\ln k}$. For $k=1$ the bound reads $0\le0$.

\emph{Part (ii).} The pairs are independent, so given $D_1,\dots,D_k$ the variable $T_i$ depends only on $D_i$, and the regression of a bivariate normal gives $\E[T_i\mid D_1,\dots,D_k]=\mu_T+\rho(\sigma_T/\sigma_D)(D_i-\mu_D)$. The index $\iota$ is a function of the $D_i$, unique with probability one, so $\E[T_\iota\mid D_1,\dots,D_k]=\mu_T+\rho(\sigma_T/\sigma_D)(\max_iD_i-\mu_D)$. Taking expectations and applying part~(i) to the $D_i$ gives $\E T_\iota=\mu_T+\rho(\sigma_T/\sigma_D)\sigma_De_k=\mu_T+\rho\sigma_Te_k$.
\end{proof}

\begin{table}[t]
\centering\small
\begin{tabular}{l*{9}{c}}
\toprule
$k$ & 2 & 3 & 4 & 5 & 6 & 7 & 8 & 9 & 10\\
\midrule
$e_k$ & 0.564 & 0.846 & 1.029 & 1.163 & 1.267 & 1.352 & 1.424 & 1.485 & 1.539\\
$\sqrt{2\ln k}$ & 1.177 & 1.482 & 1.665 & 1.794 & 1.893 & 1.973 & 2.039 & 2.096 & 2.146\\
\bottomrule
\end{tabular}
\caption{The expected maximum $e_k$ of $k$ independent standard normals (numerical integration in \script{01-seeds-bestof.py}) and the bound of \cref{prop:seeds:bestof}; further, $e_{20}=1.867$, $e_{50}=2.249$, $e_{100}=2.508$, $e_{10^4}=3.852$.}
\label{tab:seeds:ek}
\end{table}

\Cref{tab:seeds:ek} lists $e_k$, and \cref{ex:seeds:ek} derives $e_2=1/\sqrt\pi$ and $e_3=3/(2\sqrt\pi)$. The constant grows slowly, like $\sqrt{2\ln k}$ (\cref{ex:seeds:growth}), but it grows: the best of 5 runs sits $1.16\sigma$ above the mean, the best of 625 runs $3.10\sigma$. Part~(i) is the bonus of reporting the best run on the items that score it, and the bonus is luck. Part~(ii) is the honest procedure, choosing on development data and scoring once on test data: if the two scores share a real component, then the selected run is truly better, by $\rho\sigma_Te_k$, the fraction $\rho$ of the apparent gain; Reimers and Gurevych's correlation of 0.23 between development and test scores says that most of what their best development run gained was luck. The same order statistic governs the best of $k$ hyperparameter trials (\cref{ch:budgets}) and selection by a cheap proxy (\cref{ch:controls}); its asymptotics are in \cref{thm:toolkit:orderstat}.

\paragraph{Failed runs.}\index{failed run} The normal model breaks when some runs fail. If a run fails with probability $\pi$, landing near the score of a trivial predictor, then the mean mixes failures with working runs, the standard deviation is inflated by the gap between them, and the median remains the score of a working run as long as $\pi<1/2$. The best of $k$ runs fails only with probability $\pi^k$, which is why restarts worked for BERT, and why a best-of-$k$ number hides the failure rate, as the numbers for RTE in \cref{sec:seeds:postmortem} show. Report the failure rate with an exact binomial interval (\cref{thm:toolkit:clopper}), summarize the working runs separately, and show every run.

\paragraph{Variance components.} Which seeds produce the spread? A \term{crossed design}\index{crossed design} answers this: draw $I$ initialization seeds and $J$ data-order seeds, and train one run for every combination. Write $Y_{ij}$ for the outcome of the run with the $i$th initialization and the $j$th order seed, $\bar Y_{i\cdot}=\frac1J\sum_jY_{ij}$ and $\bar Y_{\cdot j}=\frac1I\sum_iY_{ij}$ for the row and column means, and $\bar Y_{\cdot\cdot}$ for the grand mean. The mean squares of the table are
\begin{equation}\label{eq:seeds:ms}
\begin{aligned}
\seedsMS A&=\frac{J}{I-1}\sum_{i=1}^I(\bar Y_{i\cdot}-\bar Y_{\cdot\cdot})^2,\qquad
\seedsMS B=\frac{I}{J-1}\sum_{j=1}^J(\bar Y_{\cdot j}-\bar Y_{\cdot\cdot})^2,\\
\seedsMS E&=\frac{1}{(I-1)(J-1)}\sum_{i=1}^I\sum_{j=1}^J(Y_{ij}-\bar Y_{i\cdot}-\bar Y_{\cdot j}+\bar Y_{\cdot\cdot})^2 .
\end{aligned}
\end{equation}
The textbook model behind them, $Y_{ij}=\mu+a_i+b_j+e_{ij}$ with independent normal effects, is an assumption. For seeds it is not needed: a deterministic training pipeline is a function of its seeds, and every such function has an exact decomposition of the same form.

\begin{lemma}[Centred sums of squares]\label{lem:seeds:centred}
If $X_1,\dots,X_n$, with $n\ge2$, are pairwise uncorrelated with a common mean and a common variance $\tau^2$, then $\E\sum_{i=1}^n(X_i-\bar X)^2=(n-1)\tau^2$.
\end{lemma}

\begin{proof}
Let $\nu$ be the common mean. Expanding around $\nu$ gives $\sum_i(X_i-\bar X)^2=\sum_i(X_i-\nu)^2-n(\bar X-\nu)^2$. The first sum has mean $n\tau^2$, and $\bar X-\nu$ has variance $\tau^2/n$ because the $X_i$ are pairwise uncorrelated, so the expectation is $n\tau^2-\tau^2$.
\end{proof}

\begin{theorem}[Variance components of a crossed design]\label{thm:seeds:components}\index{variance components}
Let $U_1,\dots,U_I$ and $V_1,\dots,V_J$, with $I,J\ge2$, be independent random seeds, the $U_i$ identically distributed and the $V_j$ identically distributed, and let $Y_{ij}=f(U_i,V_j)$ for a function $f$ with $\E f(U_1,V_1)^2<\infty$. Let $\mu=\E Y_{11}$, and define the main effects $a(u)=\E f(u,V_1)-\mu$ and $b(v)=\E f(U_1,v)-\mu$, the interaction $c(u,v)=f(u,v)-\mu-a(u)-b(v)$, and
\[
\sigma_a^2=\Var a(U_1),\qquad\sigma_b^2=\Var b(V_1),\qquad\sigma_e^2=\Var c(U_1,V_1).
\]
Then $\Var Y_{11}=\sigma_a^2+\sigma_b^2+\sigma_e^2$, and the mean squares \eqref{eq:seeds:ms} satisfy
\[
\E\seedsMS A=\sigma_e^2+J\sigma_a^2,\qquad\E\seedsMS B=\sigma_e^2+I\sigma_b^2,\qquad\E\seedsMS E=\sigma_e^2 .
\]
Consequently $\hat\sigma_a^2=(\seedsMS A-\seedsMS E)/J$, $\hat\sigma_b^2=(\seedsMS B-\seedsMS E)/I$ and $\hat\sigma_e^2=\seedsMS E$ are unbiased estimators of $\sigma_a^2$, $\sigma_b^2$ and $\sigma_e^2$.
\end{theorem}

\begin{proof}
The plan is to show that the pieces $a(U_i)$, $b(V_j)$ and $c(U_i,V_j)$ are pairwise uncorrelated, and then to compute each mean square from them with \cref{lem:seeds:centred}. Write $a_i=a(U_i)$, $b_j=b(V_j)$ and $c_{ij}=c(U_i,V_j)$, and use dots for averages over an index, as in $\bar c_{i\cdot}=\frac1J\sum_jc_{ij}$.

\emph{The pieces are orthogonal.} Because the seeds are independent, fixing $U_i=u$ leaves the law of the other seeds unchanged. Averaging $c$ over either seed gives zero: for each $u$, $\E\,c(u,V_1)=\E f(u,V_1)-\mu-a(u)-\E\,b(V_1)=0$, since $\E\,b(V_1)=\E f(U_1,V_1)-\mu=0$; symmetrically $\E\,c(U_1,v)=0$ for each $v$. Hence $a_i$, $b_j$ and $c_{ij}$ have mean zero, and they are square integrable because conditional expectations do not increase second moments. Two of them that involve disjoint sets of seeds are independent, hence uncorrelated. Otherwise the two share exactly one seed, say $U_i$ (a shared $V_j$ is symmetric), and at least one of them is a $c$. Given $U_i=u$ they are independent, and the $c$ among them, $c(u,V_j)$, has mean zero; so the conditional mean of their product is zero, and so is its expectation. In particular $\Var Y_{11}=\Var(a_1+b_1+c_{11})=\sigma_a^2+\sigma_b^2+\sigma_e^2$.

\emph{Rows and columns.} The row mean is $\bar Y_{i\cdot}=\mu+a_i+\bar b_{\cdot}+\bar c_{i\cdot}$, and $\bar b_\cdot$ is the same for every row; so $\bar Y_{i\cdot}-\bar Y_{\cdot\cdot}=X_i-\bar X$ with $X_i=a_i+\bar c_{i\cdot}$. By orthogonality the $X_i$ are pairwise uncorrelated, with mean zero and variance $\sigma_a^2+\sigma_e^2/J$, so \cref{lem:seeds:centred} gives $\E\sum_i(\bar Y_{i\cdot}-\bar Y_{\cdot\cdot})^2=(I-1)(\sigma_a^2+\sigma_e^2/J)$, and multiplying by $J/(I-1)$ gives $\E\seedsMS A=J\sigma_a^2+\sigma_e^2$. Exchanging the roles of rows and columns gives $\E\seedsMS B=I\sigma_b^2+\sigma_e^2$.

\emph{Residuals.} Substituting the decomposition into the residual, the terms in $\mu$, $a$ and $b$ cancel:
\[
r_{ij}=Y_{ij}-\bar Y_{i\cdot}-\bar Y_{\cdot j}+\bar Y_{\cdot\cdot}=c_{ij}-\bar c_{i\cdot}-\bar c_{\cdot j}+\bar c_{\cdot\cdot}.
\]
The table $c_{ij}-\bar c_{\cdot\cdot}$ is the sum of three tables, with entries $\bar c_{i\cdot}-\bar c_{\cdot\cdot}$, $\bar c_{\cdot j}-\bar c_{\cdot\cdot}$ and $r_{ij}$, which are orthogonal in $\R^{I\times J}$: the residuals sum to zero along every row and every column, the first table depends only on the row and the second only on the column, and each of these two sums to zero. By Pythagoras,
\[
\sum_{i,j}(c_{ij}-\bar c_{\cdot\cdot})^2=J\sum_i(\bar c_{i\cdot}-\bar c_{\cdot\cdot})^2+I\sum_j(\bar c_{\cdot j}-\bar c_{\cdot\cdot})^2+\sum_{i,j}r_{ij}^2 .
\]
We apply \cref{lem:seeds:centred} three times: to the $IJ$ variables $c_{ij}$, of variance $\sigma_e^2$; to the $I$ row means $\bar c_{i\cdot}$, of variance $\sigma_e^2/J$; and to the $J$ column means, of variance $\sigma_e^2/I$, all pairwise uncorrelated by orthogonality. Taking expectations,
\[
\E\sum_{i,j}r_{ij}^2=(IJ-1)\sigma_e^2-(I-1)\sigma_e^2-(J-1)\sigma_e^2=(I-1)(J-1)\sigma_e^2,
\]
which is $\E\seedsMS E=\sigma_e^2$. The three estimators are linear in the mean squares, so their unbiasedness follows.
\end{proof}

Three remarks make the theorem usable. \emph{Nothing about training is assumed} beyond independent seeds; independence and normality of the effects are needed only for $F$-tests and exact intervals (\cref{thm:toolkit:anova}). \emph{One run per cell cannot separate} the interaction from noise that is fresh in every run, such as that of nondeterministic kernels (\cref{ch:instability}); separating them takes two runs per cell. \emph{Estimates can be negative}, when $\seedsMS A<\seedsMS E$; report them as computed, since rounding them up to zero biases them. Fixing a seed leaves a conditional spread: runs that share the initialization seed spread with standard deviation $\sqrt{\sigma_b^2+\sigma_e^2}$, runs that share the order seed with $\sqrt{\sigma_a^2+\sigma_e^2}$.\sidenote{Dodge et al.\ summarized their crossed design by such conditional spreads, averaged over seeds: their ``standard deviation aggregated over weight initializations'' estimates roughly $\sqrt{\sigma_b^2+\sigma_e^2}$.} With splits as rows and runs within a split as replicates, the same analysis gives the split component.

\paragraph{How many seeds.} The last question is the one an experimenter faces first: how many runs does a comparison need?

\begin{proposition}[Power of a comparison of means]\label{prop:seeds:power}
Let each of two recipes be run $n$ times, all runs independent, with outcomes normal with means $\mu_1,\mu_2$ and a common known standard deviation $\sigma>0$, and let $\Delta=\mu_1-\mu_2$. The two-sided test that rejects $\Delta=0$ when $\abs{\bar Y_1-\bar Y_2}>z_{1-\alpha/2}\,\sigma\sqrt{2/n}$ has level $\alpha$ and power
\[
\Phi\Bigl(\frac{\sqrt n\,\abs\Delta}{\sigma\sqrt2}-z_{1-\alpha/2}\Bigr)+\Phi\Bigl(-\frac{\sqrt n\,\abs\Delta}{\sigma\sqrt2}-z_{1-\alpha/2}\Bigr).
\]
If instead the runs are paired, the $i$th runs of the two recipes sharing their seeds and split, with the pairs independent and bivariate normal and $\sigma_D^2=\Var(Y_{1i}-Y_{2i})>0$, then the same holds with $2\sigma^2$ replaced by $\sigma_D^2$, for the test that rejects when $\abs{\bar Y_1-\bar Y_2}>z_{1-\alpha/2}\,\sigma_D/\sqrt n$.
\end{proposition}

\begin{proof}
Write $s=\sigma\sqrt{2/n}$. The difference $\bar Y_1-\bar Y_2$ is normal with mean $\Delta$ and standard deviation $s$, so the rejection probability is
\[
\Prob(\bar Y_1-\bar Y_2>z_{1-\alpha/2}s)+\Prob(\bar Y_1-\bar Y_2<-z_{1-\alpha/2}s)=\Phi\Bigl(\frac\Delta s-z_{1-\alpha/2}\Bigr)+\Phi\Bigl(-\frac\Delta s-z_{1-\alpha/2}\Bigr),
\]
which is symmetric in the sign of $\Delta$ and equals the displayed power. At $\Delta=0$ it is $2\Phi(-z_{1-\alpha/2})=\alpha$. In the paired case $\bar Y_1-\bar Y_2$ is the mean of $n$ independent normal differences of variance $\sigma_D^2$, and the same computation applies with $s=\sigma_D/\sqrt n$.
\end{proof}

\begin{corollary}[Number of seeds]\label{cor:seeds:nseeds}\index{power}
In the setting of \cref{prop:seeds:power}, the power is at least $0.8$ whenever
\[
n\ge\frac{2\sigma^2\,(z_{1-\alpha/2}+z_{0.8})^2}{\Delta^2},
\]
which for $\alpha=0.05$ is $n\ge15.7\,\sigma^2/\Delta^2$ runs per arm. For paired runs, $2\sigma^2$ is replaced by $\sigma_D^2$.
\end{corollary}

\begin{proof}
The second term of the power is positive, and the first is at least $0.8$ exactly when $\sqrt n\abs\Delta/(\sigma\sqrt2)-z_{1-\alpha/2}\ge z_{0.8}$; squaring and rearranging gives the condition on $n$. With $z_{0.975}=1.960$ and $z_{0.8}=0.842$, the constant $2(z_{0.975}+z_{0.8})^2$ is $15.70$.
\end{proof}

To detect a gain of one seed standard deviation takes 16 runs per arm, of half a standard deviation 63. A gain of 0.3 points at a standard deviation of 0.2 takes 7 runs per arm by the corollary and 9 by an exact $t$-test (\cref{ex:seeds:howmany}, \cref{thm:toolkit:power}); the taggers of \cref{sec:seeds:postmortem}, with a gap of 0.27 and a standard deviation of 0.24, needed 13 each. \term{Pairing}\index{pairing} helps as much as the arms share their noise, since $\sigma_D^2=\sigma_1^2+\sigma_2^2-2\Cov(Y_1,Y_2)$: a shared split drops out of the difference, and so does the difficulty of shared test items, since only items on which two networks disagree contribute to it (\cref{thm:toolkit:mcnemar}).

\section{In the laboratory: one recipe, eighteen hundred runs}\label{sec:seeds:lab}

All three labs use one recipe: a two-layer perceptron of width 64 (8,970 parameters) on the scikit-learn digits, trained with SGD and momentum 0.9 at learning rate 0.05, batch size 64, for 20 epochs. A split puts 1,258 images in the training set and 539 in the test set, so that one misclassified digit is 0.19 points of test error.

\begin{lab}{Two hundred runs of one recipe}{01-seeds-spread}\label{lab:seeds:spread}
\labpart{Setup} Two hundred runs; run $r$ uses seed $r$ for initialization, data order and split. About 15 seconds.
\labpart{Predict first} How many test items will separate the best and the worst run? After how many runs will the 95\% interval be narrower than one test item?
\labpart{Look for} The shape of the distribution and its extremes; the running mean inside an interval that narrows like $1/\sqrt j$; the fraction of pairs of runs more than one standard deviation apart ($0.48$ by \cref{prop:seeds:tworuns}).
\end{lab}

\begin{figure}[t]
\begin{fullwidth}
\centering
\begin{tikzpicture}
\begin{groupplot}[group style={group size=2 by 1,horizontal sep=16mm},lplot,
  width=0.5\linewidth,height=0.245\linewidth]
\nextgroupplot[title={(a) two hundred runs},xlabel={test error (\%)},ylabel={number of runs},
  xmin=0.6,xmax=4.8,ymin=0,ymax=28,xtick={1,2,3,4},ytick={0,10,20}]
\addplot[reference] coordinates {(2.680,0) (2.680,28)};
\addplot[run] table[x=err,y=stack] {figures/data/01-seeds-spread-runs.dat};
\nextgroupplot[title={(b) the running mean},xmode=log,log basis x=10,xlabel={number of runs so far},
  ylabel={test error (\%)},xmin=0.85,xmax=230,ymin=0.6,ymax=4.8,minor x tick num=0,
  xtick={1,2,5,10,20,50,100,200},xticklabels={1,2,5,10,20,50,100,200},ytick={1,2,3,4}]
\addplot[name path=bandhi,draw=none,forget plot] table[x=run,y=hi] {figures/data/01-seeds-spread-runs.dat};
\addplot[name path=bandlo,draw=none,forget plot] table[x=run,y=lo] {figures/data/01-seeds-spread-runs.dat};
\addplot[ciband,forget plot] fill between[of=bandhi and bandlo];
\addplot[run] table[x=run,y=err] {figures/data/01-seeds-spread-runs.dat};
\addplot[inkblue,line width=1.0pt,no marks] table[x=run,y=mean] {figures/data/01-seeds-spread-runs.dat};
\end{groupplot}
\end{tikzpicture}
\caption{\Cref{lab:seeds:spread}: 200 runs of one recipe (two-layer MLP on the digits, 539 test images), each with new seeds for initialization, data order and split (seeds $0$--$199$). (a)~Each hollow mark is one run, runs with equal error stacked; the dashed line is the mean, $2.68\%$. (b)~The runs in the order they were made (hollow marks), the running mean (line) and its 95\% $t$-interval (band); the band for two runs, $[-1.6,7.9]$, is cut off by the frame.}
\label{fig:seeds:spread}
\end{fullwidth}
\end{figure}

\paragraph{What happened.} The 200 runs (\cref{fig:seeds:spread}a) have mean test error \pmval{2.68}{0.09} (95\% $t$-interval), standard deviation 0.65 points (95\% interval 0.60--0.73) and range 0.93\% to 4.45\%: nineteen test items separate the best and the worst run of one recipe. These are ordinary extremes; for 200 normal draws, \cref{prop:seeds:bestof} with $e_{200}=2.75$ predicts 0.88\% and 4.48\%. Of 100 disjoint pairs of runs, 50 differ by more than one standard deviation (exact 95\% interval 40--60) against 48 predicted by \cref{prop:seeds:tworuns}, with median difference 0.65 points against 0.62: two copies of this recipe, one run each, would show a gap of more than half a point about half the time.

Panel~(b) is the ornament of the chapter openings with real data: after two runs the 95\% interval is $[-1.6,7.9]$; after sixteen, all the draws of the ornament, \pmval{2.91}{0.26}; after 43 runs its half-width drops below one test item; after 200, \pmval{2.68}{0.09}, narrower than at sixteen by a factor of 2.8 rather than $\sqrt{200/16}=3.5$, because the first sixteen runs happened to spread less (standard deviation 0.48). By \cref{prop:seeds:noise}(i), a fixed network scored on a random set of 539 items measures with standard deviation 0.70 points at this error, or 0.58 for test sets drawn without replacement from the 1,797 digits: most of the spread of these runs belongs to the test set.

\begin{lab}{Where the spread comes from}{01-seeds-components}\label{lab:seeds:components}
\labpart{Setup} Crossed design: split 0, initialization seeds 0--19 crossed with data-order seeds 0--19 (400 runs). Nested design: 50 further splits, 8 runs on each with new initialization and order seeds (400 runs). Estimators of \cref{thm:seeds:components} and their one-way analogue for the split; intervals from a parametric bootstrap of the fitted normal random-effects model. About a minute.
\labpart{Predict first} Which moves the test error more in this network, the initial weights or the order of the minibatches? Will a new split move it more or less than new seeds?
\labpart{Look for} Whether the means of the initialization seeds spread out more than those of the order seeds; the interaction; the split component against the binomial noise of 539 items and its part due to item difficulty.
\end{lab}

\begin{figure}[t]
\begin{fullwidth}
\centering
\begin{tikzpicture}
\begin{groupplot}[group style={group size=2 by 1,horizontal sep=16mm},lplot,
  width=0.5\linewidth,height=0.245\linewidth,ymin=1.6,ymax=4.5,xmin=0.3,xmax=20.7,
  xtick={1,5,10,15,20},ytick={2,3,4},minor x tick num=0]
\nextgroupplot[title={(a) grouped by initialization seed},xlabel={initialization seed, by rank of its mean},ylabel={test error (\%)}]
\addplot[run] table[x=xinit,y=err] {figures/data/01-seeds-components-grid.dat};
\addplot[mean,ebars] table[x=rank,y=mean,y error=hw] {figures/data/01-seeds-components-init.dat};
\nextgroupplot[title={(b) grouped by data-order seed},xlabel={data-order seed, by rank of its mean}]
\addplot[run] table[x=xorder,y=err] {figures/data/01-seeds-components-grid.dat};
\addplot[mean,ebars] table[x=rank,y=mean,y error=hw] {figures/data/01-seeds-components-order.dat};
\end{groupplot}
\end{tikzpicture}
\caption{\Cref{lab:seeds:components}, crossed design: 400 runs on split 0, every combination of 20 initialization seeds and 20 data-order seeds. The same runs (hollow marks; equal errors side by side) grouped (a)~by initialization seed and (b)~by data-order seed, each group sorted by its mean; filled marks are group means with 95\% $t$-intervals over the 20 runs of the group.}
\label{fig:seeds:components}
\end{fullwidth}
\end{figure}

\paragraph{What happened.} The initialization seed carries 81\% of the seed variance on this split: $\hat\sigma_a=0.42$ points (95\% interval 0.29--0.55), against $\hat\sigma_b=0.04$ (0.00--0.06) for the data order and $\hat\sigma_e=0.20$ (0.18--0.21) for their interaction. The means of the initialization seeds range from 2.18\% to 3.95\%, those of the order seeds only from 2.80\% to 3.02\% (\cref{fig:seeds:components}), and the $F$-test finds the first effect beyond doubt ($p\approx10^{-125}$) and the second barely ($p=0.02$). The order matters only in combination with an initialization: runs that share the initialization seed spread with standard deviation 0.20, runs that share the order seed with 0.46. This differs from BERT fine-tuning, where \textcite{DodgeEtAl2020} found the two seeds comparable, but there the initialization seed set only the 2,048 weights of a new classifier layer, and here it sets all 8,970; the ordering is also consistent with the finding of \textcite{FrankleEtAl2020} that small networks on MNIST are stable to the noise of SGD from initialization on, an explanation we have not tested.\sidenote{Stability to SGD noise means that two runs from the same weights with different data orders end in the same linearly connected region of weights; \cref{ch:instability} takes it up.}

Is the initialization effect a property of the network or of these 539 items? On two halves of the test items, the means of the 20 initialization seeds correlate at only 0.15 (95\% interval $-0.32$ to $0.55$); \cref{lab:seeds:bestof} settles the question. A new split moves the error about as much as new seeds do: 0.44 points (0.35--0.53) between splits against 0.38 (0.35--0.41) within them. The binomial noise of 539 items is 0.73 points at this error (0.61 for the finite pool), and by \cref{prop:seeds:noise}(iii) its part that a new split moves and new seeds do not, item difficulty, is 0.52 points after the finite-pool correction, inside the interval for the split: a new split acts almost entirely through which digits land in the test set, as \textcite{BouthillierEtAl2021} found for data sampling at larger scale.

\begin{lab}{The best of $k$}{01-seeds-bestof}\label{lab:seeds:bestof}
\labpart{Setup} Split 0, 400 runs with new initialization and order seeds, once on clean labels and once with 20\% of the training labels replaced by random labels. The 539 test images are divided at random into a development half (270) and a test half (269). For each $k\le50$, exact averages over all $k$-subsets of the 400 runs of the best test-half error, and of the test-half error of the run with the best development-half error (ties broken at random). About 70 seconds.
\labpart{Predict first} On clean labels, how much lower is the best of 10 runs than a typical run? How much of that survives when the best run is chosen on the development half and scored on the test half? And with noisy labels?
\labpart{Look for} The best-of-$k$ curve against $\mu-\sigma e_k$, and where it departs; the honest curve against $\mu-\rho\sigma e_k$; the correlation $\rho$ between development and test errors in the two settings.
\end{lab}

\begin{figure}[t]
\begin{fullwidth}
\centering
\begin{tikzpicture}
\begin{groupplot}[group style={group size=2 by 1,horizontal sep=16mm},lplot,
  width=0.5\linewidth,height=0.245\linewidth,xmode=log,log basis x=10,xmin=0.85,xmax=60,
  minor x tick num=0,xtick={1,2,5,10,20,50},xticklabels={1,2,5,10,20,50},
  xlabel={number of runs $k$}]
\nextgroupplot[title={(a) clean labels},ylabel={test-half error (\%)},ymin=1.6,ymax=3.4,ytick={2,2.5,3},
  yticklabels={2.0,2.5,3.0},legend to name=seedsbestoflegend,legend columns=2,
  legend style={/tikz/every even column/.append style={column sep=5mm}}]
\addplot[name path=bh,draw=none,forget plot] table[x=k,y=bhi] {figures/data/01-seeds-bestof-clean.dat};
\addplot[name path=bl,draw=none,forget plot] table[x=k,y=blo] {figures/data/01-seeds-bestof-clean.dat};
\addplot[fill=rubric!12,draw=none,forget plot] fill between[of=bh and bl];
\addplot[name path=hh,draw=none,forget plot] table[x=k,y=hhi] {figures/data/01-seeds-bestof-clean.dat};
\addplot[name path=hl,draw=none,forget plot] table[x=k,y=hlo] {figures/data/01-seeds-bestof-clean.dat};
\addplot[fill=inkblue!12,draw=none,forget plot] fill between[of=hh and hl];
\addplot[rubric,line width=0.8pt,no marks] table[x=k,y=pred] {figures/data/01-seeds-bestof-clean.dat};
\addlegendentry{$\mu-\sigma e_k$}
\addplot[only marks,mark=square*,mark size=1.7pt,color=rubric,mark options={fill=rubric,draw=white,line width=0.2pt}] table[x=k,y=best] {figures/data/01-seeds-bestof-clean.dat};
\addlegendentry{best of $k$, scored on the items that chose it}
\addplot[inkblue,line width=0.8pt,dashed,no marks] table[x=k,y=hpred] {figures/data/01-seeds-bestof-clean.dat};
\addlegendentry{$\mu-\rho\sigma e_k$}
\addplot[only marks,mark=*,mark size=1.8pt,color=inkblue,mark options={fill=inkblue,draw=white,line width=0.2pt}] table[x=k,y=honest] {figures/data/01-seeds-bestof-clean.dat};
\addlegendentry{best of $k$ on development items, scored on test items}
\nextgroupplot[title={(b) 20\% of training labels random},ymin=9.5,ymax=15.5,ytick={10,12,14}]
\addplot[name path=bh,draw=none,forget plot] table[x=k,y=bhi] {figures/data/01-seeds-bestof-noisy.dat};
\addplot[name path=bl,draw=none,forget plot] table[x=k,y=blo] {figures/data/01-seeds-bestof-noisy.dat};
\addplot[fill=rubric!12,draw=none,forget plot] fill between[of=bh and bl];
\addplot[name path=hh,draw=none,forget plot] table[x=k,y=hhi] {figures/data/01-seeds-bestof-noisy.dat};
\addplot[name path=hl,draw=none,forget plot] table[x=k,y=hlo] {figures/data/01-seeds-bestof-noisy.dat};
\addplot[fill=inkblue!12,draw=none,forget plot] fill between[of=hh and hl];
\addplot[rubric,line width=0.8pt,no marks] table[x=k,y=pred] {figures/data/01-seeds-bestof-noisy.dat};
\addplot[only marks,mark=square*,mark size=1.7pt,color=rubric,mark options={fill=rubric,draw=white,line width=0.2pt}] table[x=k,y=best] {figures/data/01-seeds-bestof-noisy.dat};
\addplot[inkblue,line width=0.8pt,dashed,no marks] table[x=k,y=hpred] {figures/data/01-seeds-bestof-noisy.dat};
\addplot[only marks,mark=*,mark size=1.8pt,color=inkblue,mark options={fill=inkblue,draw=white,line width=0.2pt}] table[x=k,y=honest] {figures/data/01-seeds-bestof-noisy.dat};
\end{groupplot}
\end{tikzpicture}

\smallskip
\pgfplotslegendfromname{seedsbestoflegend}
\caption{\Cref{lab:seeds:bestof}: the best of $k$ runs, 400 runs per panel (split 0, seeds $0$--$399$), errors on a test half of 269 images. Squares: expected lowest test-half error among $k$ runs; circles: expected test-half error of the run with the lowest error on the other 270 images; both exact averages over all $k$-subsets, with 95\% bootstrap bands over runs. Lines: the predictions of \cref{prop:seeds:bestof} from the sample mean $\mu$, standard deviation $\sigma$ and development--test correlation $\rho$ of each panel.}
\label{fig:seeds:bestof}
\end{fullwidth}
\end{figure}

\paragraph{What happened.} On clean labels (\cref{fig:seeds:bestof}a) the test-half error has mean \pmval{3.06}{0.06} and standard deviation 0.56 points. The best of $k$ runs scored on the test half follows $\mu-\sigma e_k$: for $k=10$ it is 2.24\% (95\% band 2.18--2.32) against 2.19\% predicted, 0.82 points below a typical run; at $k=50$ it falls short of the prediction (1.94\% against 1.79\%), because the errors are skewed to the right (skewness 0.23) and cannot fall below the few digits every run gets wrong. Chosen on the development half, the best of 10 runs scores 3.05\% (2.92--3.19) on the test half: nothing of the 0.82 points survives, because a run's errors on the two halves are uncorrelated ($\rho=-0.03$, 95\% interval $-0.13$ to $0.07$). Item churn makes up 91\% of the seed variance on all 539 items (80--104\%): on clean digits the seed decides which borderline digits a network gets wrong, not how many it would get wrong on new digits, and the initialization effect of \cref{lab:seeds:components} is an effect on which items fail.

With 20\% of the training labels random (\cref{fig:seeds:bestof}b) the runs differ for real: mean \pmval{15.06}{0.22}, standard deviation 2.20 points, $\rho=0.51$ (0.43--0.58), and item churn only 37\% of the seed variance (32--43\%). The best of 10 runs scores 11.86\% on the items that chose it (11.69\% predicted) and 13.58\% (13.19--13.97) when chosen on the development half (13.35\% predicted by $\mu-\rho\sigma e_k$). At every $k$ the honest gain is 44--48\% of the apparent gain, close to $\rho$. Restarts selected on held-out data buy a real improvement when runs differ in what they learned, as they did for BERT; scored on the items that selected them, they report that improvement plus a bonus of luck.

The miniature reproduced the post-mortems, and differed from them in one respect: here the initialization seed outweighs the data-order seed.

\section{Protocol}\label{sec:seeds:protocol}

\begin{protocol}{Measuring one recipe}\label{prot:seeds:measure}
\item Decide which sources of randomness the claim must hold across (initialization, data order, augmentation and dropout, split, hardware); randomize those, fix the rest, and record every seed.
\item Fix the number of runs before looking, from a pilot standard deviation and the smallest difference worth detecting (\cref{cor:seeds:nseeds}); never fewer than five for a description.
\item Report the mean with a 95\% interval, the standard deviation, the number of runs and every run, in a dot plot or a table; report failed runs and the rule that defines them separately.
\item Never report the best of $k$ runs without $k$ and the mean; when a run must be selected, select it on development data and score it once on test data.
\item Compare recipes on the same splits, seeds and test items, and test the difference across runs, not only across test items.
\item Compute the binomial noise of the test set and, from per-item predictions, its churn; if they are comparable to the effect under study, then the test set is too small for the claim.
\item When the spread matters, estimate where it comes from with a crossed design before choosing which seeds to vary.
\end{protocol}

The protocol makes the measurement of one recipe honest; it does not make a comparison fair. Two well-measured recipes can still differ in their tuning budgets (\cref{ch:budgets}) or in more than the component a claim is about (\cref{ch:confounding}), be compared on many tasks at once (\cref{ch:aggregates}), or behave differently on other hardware (\cref{ch:instability}). And at 16 runs per arm to detect a gain of one standard deviation, the protocol is out of reach for the largest models. Then measure the spread where runs are affordable (a smaller model, fewer steps, a subset of the data) and say so, use paired designs, and when only one run is possible, report it as one run beside the seed standard deviation of a comparable setting, with a provisional conclusion.

\begin{exercises}
\item \consolidation (Expected maxima.) (a) Show that $e_2=1/\sqrt\pi$, using $\max(a,b)=\frac12(a+b)+\frac12\abs{a-b}$, and that $e_3=3/(2\sqrt\pi)$. (b) Show that $e_k=\int zk\varphi(z)\Phi(z)^{k-1}\dd z$, and check \cref{tab:seeds:ek} numerically. (c) Show that the expected range of $k$ independent $\cN(\mu,\sigma^2)$ draws is $2\sigma e_k$, and check the CIFAR-10 range of \textcite{Picard2021}. (d) For runs uniform on $[\mu-h,\mu+h]$, show that the best of $k$ has mean $\mu+h(k-1)/(k+1)$, and compare its growth with that of $e_k$; which run distributions make best-of-$k$ reports most misleading? \hint{for $e_3$, integrate $3\int z\varphi\Phi^2\dd z$ by parts, then use $\Phi(z)+\Phi(-z)=1$ and $\int\varphi^2=1/(2\sqrt\pi)$.}\label{ex:seeds:ek}

\item \consolidation\hard (The growth of $e_k$.) Prove that $e_k/\sqrt{2\ln k}\to1$. \hint{for the lower bound, with $M=\max_iZ_i$ and $t=(1-\eps)\sqrt{2\ln k}$, show $e_k\ge t\Prob(M\ge t)-\E\abs{Z_1}$, and use $1-\Phi(t)\ge\varphi(t)\,t/(1+t^2)$ to show $\Phi(t)^k\to0$.}\label{ex:seeds:growth}

\item \consolidation (How many seeds.) A method is expected to gain 0.3 points over a baseline whose seed standard deviation is 0.2 points. (a) By \cref{cor:seeds:nseeds}, how many runs per arm give power 0.8 at level 0.05? (b) With the noncentral $t$ distribution, show that a two-sample $t$-test needs 9 runs per arm. (c) Suppose the runs are paired by seed and split, and paired outcomes have correlation 0.6. How many pairs suffice by the corollary?\label{ex:seeds:howmany}

\item \labwork (A smaller test set.) Modify \script{01-seeds-spread.py} to score every run on only the first 100 test images of its split. Predict the new standard deviation from \cref{prop:seeds:noise}(i) and \cref{lab:seeds:spread} before running it, then compare.\label{ex:seeds:smalltest}

\item \labwork (When does the order matter?) Rerun the crossed design of \script{01-seeds-components.py} with learning rate 0.2, then with batch size 16, then with Adam at learning rate $10^{-3}$. Observe how $\hat\sigma_b$ and $\hat\sigma_e$ change relative to $\hat\sigma_a$, and relate the change to the stability of training to the noise of SGD (\cref{ch:instability}).\label{ex:seeds:order}

\item \labwork (The size of the development set.) In \script{01-seeds-bestof.py} with noisy labels, use development sets of $d=50$, $270$ and $489$ of the 539 test images and test sets of the rest. Predict how $\rho$ changes with $d$, check it, and observe the honest best-of-10 error. When is selecting the best of $k$ runs on development data worth its cost?\label{ex:seeds:devsize}

\item \design (Components of a fine-tuning pipeline.) A team fine-tunes a pretrained language model and can vary the pretraining checkpoint (five are available, differing only in their pretraining seed), the initialization of the new head and the data order, with a budget of 300 runs. Design a study of the three main components: which interactions are estimable and which are confounded with what, how the runs are allocated, and how precisely five checkpoints determine their component.\label{ex:seeds:design}

\item \design (Paired or unpaired.) Two recipes are to be compared with 60 runs in all. A pilot gives a seed standard deviation of 0.5 points for each and a correlation of 0.7 between runs of the two recipes that share split, seeds and test items. Compare the power of paired and unpaired designs to detect a gain of 0.2 points, and list what the code must do for the pairing to be real.\label{ex:seeds:pair}

\item \critique (Best of five against single runs.) A paper's main table lists its method as ``best of 5 seeds'' next to baseline scores copied from the original papers, each a single run. The method leads by 0.4 points, and its own runs have a standard deviation of 0.3. What does the lead measure? Estimate the part due to selection with \cref{prop:seeds:bestof}, and describe the table a reader needs.\label{ex:seeds:critique1}

\item \critique (The best and the worst seed.) A study crosses 25 initialization seeds with 25 data orders, picks the initialization seeds with the highest and the lowest mean, and reports a two-sample $t$-test between their 25 runs each, $p<0.01$, as evidence that some initializations are better than others. Suppose that all runs are independent with standard deviation $\sigma$, as when neither seed has an effect. Show that the expected gap between the best and the worst seed mean is $2e_{25}\sigma/5\approx0.79\sigma$ ($e_{25}=1.965$), while the standard error of the test is $\sigma\sqrt{2/25}\approx0.28\sigma$. Which test answers the question?\label{ex:seeds:critique2}

\item \research (Intervals for bimodal run distributions.) When some runs fail, the mean is a poor summary and the $t$-interval loses its coverage. Open: a procedure that reports the failure rate and the mean of working runs jointly, with a failure rule that does not look at test data and with guaranteed coverage when failures are rare, at rates near $1/n$. Also open: a theory that predicts which seed dominates the spread (the initialization in \cref{lab:seeds:components}, neither for BERT fine-tuning) from the width, the learning rate and the fraction of weights a seed initializes.\label{ex:seeds:bimodal}
\end{exercises}

\notesandrefs
The problem of this chapter is older than deep learning. \Textcite{Student1908} derived the $t$ distribution for the mean of a few measurements, and \textcite{Tippett1925} studied the extremes and the range of normal samples, whose mean is $2e_k$ in our notation. \Cref{thm:seeds:components} is the classical result on expected mean squares in the balanced two-way layout \parencite{SearleEtAl1992}, stated here for a function of independent seeds.

For deep networks, \textcite{ReimersGurevych2017} put the seed on the agenda of natural language processing, and \textcite{CardEtAl2020} showed that many comparisons in the field are underpowered; in reinforcement learning, \textcite{HendersonEtAl2018} (\cref{ch:prologue}) and \textcite{ColasEtAl2018} raised the same questions, the latter with a power analysis for the number of seeds. \Textcite{BouthillierEtAl2019} distinguished reproducing a result from reproducing a finding, and \textcite{GundersenEtAl2022} review the sources of irreproducibility in machine learning. Seeds change more than scores: they change the explanations of predictions \parencite{MadhyasthaJain2019} and generalization outside the training distribution \parencite{McCoyEtAl2020}. That a best-of-$k$ number depends on $k$ is the starting point of \textcite{DodgeEtAl2019} and of \cref{ch:budgets}.

\paragraph{Translating the literature.} \Textcite{ReimersGurevych2017} report F1 as a fraction, so their $\sigma=0.00241$ is our 0.24 points; their score distribution is our run distribution, and the ``Evaluation~1'' and ``Evaluation~2'' of their 2018 preprint are our one run per arm and best of $k$ on development data. \Textcite{DodgeEtAl2020} call a run a trial, while in this book a trial is a hyperparameter configuration; their WI and DO seeds are our initialization and order seeds, their expected validation performance of the best of $x$ trials is our $\E[M_x]$ on the validation set, and their standard deviation aggregated over WI seeds is the conditional spread $\sqrt{\sigma_b^2+\sigma_e^2}$ of \cref{thm:seeds:components}, averaged over seeds. What \textcite{Picard2021} calls the seed is our full set of seeds. In \textcite{BouthillierEtAl2021}, data sampling is our split component, $\xi_O$ collects our initialization, order and noise seeds, and numerical noise is the nondeterminism of \cref{ch:instability}.

\printchapterbib
\end{refsection}
